\section{Аналитическая часть}

В данной части проводится анализ объектов сцены и существующих алгоритмов построения изображений и выбор более подходящих алгоритмов для дальнейшего использования.

\subsection{Формализация объектов синтезируемой сцены}

Сцена театрального оформления состоит из определённого набора объектов, каждый из которых обладает геометрическими и спектральными характеристиками.

\begin{enumerate}
	\item \textbf{Сценическая площадка} --- базовая плоскость (или система плоскостей), задающая геометрические границы сцены: глубину, ширину, высоту порталов и рабочую зону. Площадка может содержать дополнительные элементы (подиумы, люки, рельсы), однако в рамках проекта достаточно аппроксимировать её набором плоских граней.
	
	\item \textbf{Декорационные элементы} --- объекты, формирующие художественное оформление сцены. К ним относятся:
	\begin{itemize}
		\item жёсткие декорации (стенки, объёмные конструкции), которые удобно представлять многогранниками или полигональными сетками;
		\item мягкие декорации (занавесы, драпировки, задники), которые в первом приближении также аппроксимируются полигональными поверхностями;
		\item реквизит и мебель (столы, стулья, колонны и т.\,п.), задаваемые параметрически или в виде готовых полигональных моделей.
	\end{itemize}
	Каждый декорационный элемент характеризуется геометрическими параметрами (размеры, положение, ориентация) и спектральными характеристиками (диффузный и зеркальный коэффициенты отражения, цвет, коэффициент блеска).
	
	\item \textbf{Источники света} --- театральные осветительные приборы. В проекте рассматриваются:
	\begin{itemize}
		\item направленные источники (прожекторы), задаваемые положением, направлением, углом раскрытия луча и RGB-параметрами цвета;
		\item точечные источники;
		\item рассеянный свет, моделирующий общее освещение сцены.
	\end{itemize}
	От расположения и параметров источников зависят тени и общая световая картина.
	
	\item \textbf{Камера} --- фиксированный виртуальный наблюдатель с заранее заданными положением, направлением взгляда и параметрами проекции (угол обзора, ближняя и дальняя плоскости отсечения). Перемещение и вращение камеры пользователем не предусмотрено.
\end{enumerate}

Таким образом, синтезируемая сцена представляет собой композицию из полигональных объектов, расположенных относительно сценической площадки, с набором управляемых источников света и камеры.

\subsection{Способы описания трёхмерных геометрических моделей на сцене}

В компьютерной графике для описания трёхмерных геометрических объектов применяют три основных типа моделей: каркасную, поверхностную и твердотельную.

\begin{itemize}
	\item \textbf{Каркасная модель} содержит минимум информации --- только о вершинах и рёбрах. Она экономична по памяти и вычислительным затратам, однако не позволяет корректно определять видимость граней, различать внутреннюю и внешнюю стороны поверхности, а также корректно строить закраску и тени. Для задач моделирования театральных декораций, где важны поверхности и освещение, каркасная модель неприменима.
	
	\item \textbf{Поверхностная модель} дополняет каркасную информацией о гранях (полигонах), ограничивающих объект. Она широко используется в интерактивной компьютерной графике, поскольку позволяет реализовать удаление невидимых поверхностей, закраску и расчёт освещения. Недостатком является отсутствие явного указания, с какой стороны поверхности находится материал (хотя направление нормали обычно решает эту проблему).
	
	\item \textbf{Твердотельная модель} дополнительно хранит информацию о том, какая часть пространства принадлежит материалу объекта (обычно через направление внутренней нормали или конструктивную твердотельную геометрию). Для задач визуализации оформления сцены такая информация избыточна, поскольку не требуется моделирование физических свойств материалов или булевых операций над объёмами.
\end{itemize}

Для решения поставленной задачи наиболее подходящей является \textbf{поверхностная модель}, так как она обеспечивает необходимый баланс между информативностью и вычислительной сложностью, позволяя корректно отображать форму декораций, рассчитывать освещение и тени.

\subsection{Способы задания поверхностных моделей}

Поверхностная модель может быть задана аналитическим (параметрическим) способом либо в виде полигональной сетки.

\begin{itemize}
	\item \textbf{Аналитический способ} описывает поверхность неявными или параметрическими уравнениями. Он удобен для поверхностей вращения и гладких форм, однако при интерактивном редактировании и растеризации требует дополнительной аппроксимации полигонами.
	
	\item \textbf{Полигональная сетка} представляет объект совокупностью вершин, рёбер и граней. Это основной способ представления в системах реального времени.
\end{itemize}

Для хранения полигональной сетки существуют различные схемы:

\begin{itemize}
	\item вершинное представление (список вершин и граней, ссылающихся на индексы вершин);
	\item рёберное представление (winged-edge и его варианты);
	\item представление в виде списка граней с явным хранением координат.
\end{itemize}

В рамках проекта целесообразно использовать представление в виде списка вершин и списка граней (индексированная полигональная сетка). Такое представление компактно, удобно для аффинных преобразований и хорошо сочетается с алгоритмом Z-буфера.

Декорационные элементы могут создаваться параметрически (например, прямоугольные стенки, цилиндрические колонны, призматические конструкции) с последующим преобразованием в полигональную сетку, либо загружаться из внешних файлов.

\subsection{Анализ алгоритмов удаления невидимых линий и поверхностей}

Для корректного отображения трёхмерной сцены необходимо удалять невидимые линии и поверхности. Рассмотрим наиболее известные алгоритмы.

\subsubsection{Алгоритм Робертса}

Алгоритм Робертса работает в объектном пространстве и предназначен преимущественно для выпуклых многогранников. Он анализирует рёбра и определяет их видимость относительно других граней. Алгоритм требует предварительной проверки выпуклости объектов и плохо масштабируется при большом числе пересекающихся объектов. Для интерактивного редактирования сцены с произвольными декорациями он малопригоден.

\subsubsection{Алгоритм Варнока}

Алгоритм Варнока работает в пространстве изображения. Экран рекурсивно разбивается на области до тех пор, пока в каждой области не останется один из простых случаев: пусто или один полигон. Алгоритм эффективен для сцен с большими однородными областями, однако рекурсивное деление и сложность обработки пересечений делают его менее предсказуемым по времени выполнения.

\subsubsection{Алгоритм, использующий Z-буфер}

Алгоритм Z-буфера работает в пространстве изображения. Для каждого пикселя хранится значение глубины ближайшего к наблюдателю фрагмента. При растеризации полигона для каждого пикселя сравнивается текущее значение глубины с значением в буфере; при необходимости буфер обновляется. Алгоритм прост в реализации, хорошо параллелизуется, его трудоёмкость слабо зависит от числа объектов (в основном от разрешения экрана и площади проекций). Недостатки --- потребность в памяти под Z-буфер и возможное появление артефактов при недостаточной точности хранения глубины. Для интерактивных приложений моделирования сцены алгоритм Z-буфера является одним из наиболее подходящих.

\subsubsection{Алгоритм прямой трассировки лучей}

При прямой трассировке лучи испускаются от источников света и отслеживаются до попадания в камеру. Алгоритм способен моделировать сложные оптические эффекты, однако крайне неэффективен, поскольку большинство лучей не попадает в камеру. Для интерактивного моделирования театральной сцены неприменим.

\subsubsection{Алгоритм обратной трассировки лучей}

Обратная трассировка лучей испускает первичные лучи из камеры через каждый пиксель и ищет ближайшее пересечение с объектами сцены. Для расчёта освещения и теней испускаются вторичные лучи к источникам света. Алгоритм даёт высокое качество изображения и естественным образом обрабатывает отражения, преломления и тени, однако вычислительно затратен и плохо подходит для интерактивного редактирования при ограниченных ресурсах.

\subsubsection{Выбор оптимального алгоритма}

Учитывая требования интерактивности (перемещение декораций и источников света), ограничения по сложности реализации и необходимость работы с произвольным числом полигональных объектов, оптимальным выбором является \textbf{алгоритм Z-буфера}. Он обеспечивает приемлемое качество визуализации при умеренных вычислительных затратах и простоте программной реализации.

\subsection{Анализ алгоритмов закраски}

После определения видимых поверхностей необходимо выполнить закраску.

\subsubsection{Простая закраска}

Интенсивность освещения вычисляется один раз для всей грани (обычно по нормали в центре). Метод вычислительно простой, однако подчёркивает полигональную структуру модели и не обеспечивает плавных переходов.

\subsubsection{Закраска по Гуро}

Интенсивность вычисляется в вершинах грани, а затем линейно интерполируется вдоль рёбер и по строкам развёртки. Метод даёт более гладкое изображение, чем плоская закраска, хорошо маскирует полигональность, но может пропускать блики внутри грани и даёт менее точные блики.

\subsubsection{Закраска по Фонгу}

Интерполируются нормали (а не интенсивности) по поверхности грани, после чего в каждой точке вычисляется освещённость. Метод обеспечивает более корректные блики и более высокое качество, но требует больших вычислительных затрат.

\subsubsection{Выбор алгоритма закраски}

Для программного комплекса моделирования театральной сцены, где важны и интерактивность, и приемлемое качество восприятия света на декорациях, целесообразно использовать \textbf{закраску по Гуро} как компромисс между качеством и скоростью.

\subsection{Анализ алгоритмов построения теней}

Тени существенно повышают реалистичность восприятия сцены и помогают оценить взаимное расположение декораций и света. Основные подходы:

\begin{itemize}
	\item теневые объёмы;
	\item карты теней;
	\item построение теней в процессе трассировки лучей.
\end{itemize}

Для алгоритма Z-буфера наиболее естественным и относительно простым в реализации является метод \textbf{карт теней}: сцена рендерится из положения источника света в буфер глубины, после чего при основном рендеринге для каждого фрагмента проверяется, находится ли он в тени. Метод хорошо сочетается с Z-буфером и направленными/точечными источниками, характерными для театрального освещения.

\subsection{Анализ моделей освещения}

Модели локального освещения определяют, как свет взаимодействует с поверхностью.

\subsubsection{Модель Ламберта}

Диффузное отражение пропорционально косинусу угла между нормалью и направлением на источник света. Модель проста и хорошо описывает матовые поверхности, но не учитывает зеркальные блики.

\subsubsection{Модель Фонга}

Дополняет диффузную составляющую зеркальной: интенсивность блика зависит от угла между отражённым лучом и направлением на наблюдателя, возведённым в степень блеска. Модель широко используется в интерактивной графике и позволяет имитировать блестящие поверхности декораций и металлические элементы.

\subsubsection{Выбор модели освещения}

В проекте принимается \textbf{модель освещения Фонга} (с диффузной и зеркальной составляющими, а также ambient-членом, отвечающим за фоновое освещение). Она обеспечивает достаточную выразительность световой картины театральной сцены при умеренной вычислительной сложности.

\subsection{Постановка задачи}

На основании проведённого анализа формулируется постановка задачи.

\textbf{Входные данные:}
\begin{itemize}
	\item параметры сценической площадки;
	\item набор декорационных элементов с геометрическими и спектральными характеристиками;
	\item параметры источников света (положение, направление, цвет, интенсивность);
	\item параметры камеры (заданные при инициализации и неизменяемые).
\end{itemize}

\textbf{Выходные данные:}
\begin{itemize}
	\item растровое изображение сцены с учётом удаления невидимых поверхностей, закраски и теней;
	\item возможность интерактивного изменения параметров объектов декораций и света.
\end{itemize}

\textbf{Ограничения:}
\begin{itemize}
	\item объекты представляются полигональными сетками;
	\item число объектов и полигонов ограничено требованиями интерактивности;
	\item используется локальная модель освещения (без глобального освещения и сложных оптических эффектов).
\end{itemize}

\textbf{Критерии сравнения реализаций:}
\begin{itemize}
	\item время построения кадра в зависимости от числа объектов и полигонов;
	\item визуальное качество (корректность теней, плавность закраски).
\end{itemize}

\textbf{Функциональные требования к программному обеспечению:}
\begin{itemize}
	\item создание, редактирование и удаление декорационных элементов;
	\item управление положением и параметрами источников света;
	\item визуализация сцены с использованием выбранных алгоритмов;
	\item сохранение и загрузка сцены.
\end{itemize}

\newpage